\chapter*{Further reading}\addcontentsline{toc}{chapter}{Further reading}\markboth{Further reading}{}
For the geometry of local fields, start with the early chapters of Giachetta, Mangiarotti and Sardanashvily~\cite{GMS}: the reason to learn a jet bundle is to give a coordinate-independent domain to a differential equation or Lagrangian. Lee~\cite{Lee} and Steenrod~\cite{Steenrod} supply background on smooth manifolds and bundle clutching. At a null face a Lorentzian reference such as O'Neill~\cite{ON} is useful precisely where Riemannian intuitions about normal vectors fail.

For boundary variations the convention used here follows Harlow and Wu~\cite{HW}. Khavkine~\cite{Khavkine} explains the dual descriptions of brackets and the conditions needed in the presence of constraints. Cattaneo, Mnev and Reshetikhin~\cite{CMR} explain why a boundary in BV theory has its own structure, while Donnelly and Freidel~\cite{DF} develop an extended description of subsystems. Neither reference licenses treating an arbitrary boundary gauge transformation as a proper gauge degeneracy. For wave operators, B\"ar, Ginoux and Pf\"affle~\cite{BGP} treat the globally hyperbolic Cauchy problem; B\"ar and Tagne Wafo~\cite{BW} treat characteristic data under additional geometric and support hypotheses. Neither is a universal timelike transmission theorem.

For the distinction between a field algebra and a choice of state, compare Wald~\cite{WaldQFT} with Haag~\cite{Haag}. After the elementary scaling calculation in Chapter~\ref{ch:13}, Brunetti and Fredenhagen~\cite{BF} and Hollands and Wald~\cite{HWick,HWT} are natural points of entry into causal renormalization and its local freedom. Their results on smooth curved spacetimes are not existence theorems for all possible physical boundaries, corners or defects.

\renewcommand{\bibname}{References}
\begin{thebibliography}{99}
\addcontentsline{toc}{chapter}{References}
\bibitem{GMS} G. Giachetta, L. Mangiarotti and G. Sardanashvily, \emph{Advanced Classical Field Theory}, World Scientific (2009).
\bibitem{Lee} J. M. Lee, \emph{Introduction to Smooth Manifolds}, 2nd ed., Springer (2013).
\bibitem{Steenrod} N. Steenrod, \emph{The Topology of Fibre Bundles}, Princeton University Press (1951).
\bibitem{ON} B. O'Neill, \emph{Semi-Riemannian Geometry with Applications to Relativity}, Academic Press (1983).
\bibitem{HW} D. Harlow and J.-q. Wu, ``Covariant phase space with boundaries,'' JHEP \textbf{10} (2020) 146, \href{https://arxiv.org/abs/1906.08616}{arXiv:1906.08616}.
\bibitem{Khavkine} I. Khavkine, ``Covariant phase space, constraints, gauge and the Peierls formula,'' Int. J. Mod. Phys. A \textbf{29} (2014) 1430009, \href{https://arxiv.org/abs/1402.1282}{arXiv:1402.1282}.
\bibitem{CMR} A. S. Cattaneo, P. Mnev and N. Reshetikhin, ``Classical BV theories on manifolds with boundary,'' Commun. Math. Phys. \textbf{332} (2014) 535--603, \href{https://arxiv.org/abs/1201.0290}{arXiv:1201.0290}.
\bibitem{DF} W. Donnelly and L. Freidel, ``Local subsystems in gauge theory and gravity,'' JHEP \textbf{09} (2016) 102, \href{https://arxiv.org/abs/1601.04744}{arXiv:1601.04744}.
\bibitem{BGP} C. B\"ar, N. Ginoux and F. Pf\"affle, \emph{Wave Equations on Lorentzian Manifolds and Quantization}, European Mathematical Society (2007).
\bibitem{BW} C. B\"ar and R. Tagne Wafo, ``Initial value problems for wave equations on manifolds,'' \emph{Math. Phys. Anal. Geom.} \textbf{18} (2015), article 7, \href{https://arxiv.org/abs/1408.4995}{arXiv:1408.4995}.
\bibitem{WaldQFT} R. M. Wald, \emph{Quantum Field Theory in Curved Spacetime and Black Hole Thermodynamics}, University of Chicago Press (1994).
\bibitem{Haag} R. Haag, \emph{Local Quantum Physics: Fields, Particles, Algebras}, 2nd ed., Springer (1996).
\bibitem{Rad} M. J. Radzikowski, ``Micro-local approach to the Hadamard condition in quantum field theory on curved space-time,'' Commun. Math. Phys. \textbf{179} (1996) 529--553.
\bibitem{BF} R. Brunetti and K. Fredenhagen, ``Microlocal analysis and interacting quantum field theories: renormalization on physical backgrounds,'' Commun. Math. Phys. \textbf{208} (2000) 623--661, \href{https://arxiv.org/abs/math-ph/9903028}{arXiv:math-ph/9903028}.
\bibitem{HWick} S. Hollands and R. M. Wald, ``Local Wick polynomials and time ordered products of quantum fields in curved spacetime,'' Commun. Math. Phys. \textbf{223} (2001) 289--326, \href{https://arxiv.org/abs/gr-qc/0103074}{arXiv:gr-qc/0103074}.
\bibitem{HWT} S. Hollands and R. M. Wald, ``Existence of local covariant time ordered products of quantum fields in curved spacetime,'' Commun. Math. Phys. \textbf{231} (2002) 309--345, \href{https://arxiv.org/abs/gr-qc/0111108}{arXiv:gr-qc/0111108}.
\bibitem{FR} K. Fredenhagen and K. Rejzner, ``Perturbative algebraic quantum field theory,'' \href{https://arxiv.org/abs/1208.1428}{arXiv:1208.1428}; see also K. Rejzner, \emph{Perturbative Algebraic Quantum Field Theory}, Springer (2016).
\end{thebibliography}
